\section{Problem Formulation and Planning Background}~\label{preliminaries}

\subsection{Symbolic Task Model}~\label{preliminaries:STRIPS}
A symbolic planning problem is a tuple $\mathcal{P}=\langle\mathcal{F},\mathcal{A},s_0,\mathcal{G}\rangle$, where $\mathcal{F}$ is a finite set of atoms, $\mathcal{A}$ is a set of actions, $s_0\subseteq\mathcal{F}$ is the initial state, and $\mathcal{G}\subseteq\mathcal{F}$ is the goal condition. A state $s\subseteq\mathcal{F}$ specifies the atoms that hold. Each action $a\in\mathcal{A}$ is defined by its preconditions and effects,
\[
a=\langle\mathrm{pre}(a),\mathrm{add}(a),\mathrm{del}(a)\rangle.
\]
The action is applicable when $\mathrm{pre}(a)\subseteq s$, denoted by $a\llbracket s\rrbracket$, and produces
\[
s'=(s\setminus\mathrm{del}(a))\cup\mathrm{add}(a).
\]
A plan is an action sequence that transforms $s_0$ into a state satisfying $\mathcal{G}$.

Our setting departs from classical planning because the current symbolic state is not fully known. After executing an action, multiple successor states may remain plausible because of stochastic action outcomes and noisy observations. The robot must therefore make two feedback decisions during execution: whether the remaining state ambiguity warrants an external query and, if so, which predicate fact should be verified. Queries return binary evidence about the current world state; they do not directly specify the robot's next action.

\subsection{Belief-Space Online Planning}~\label{preliminaries:POMDPs}
We model sequential decision making with a POMDP $\langle\mathcal{S},\mathcal{A},T,R,\Omega,O,\gamma\rangle$. The transition model $T(s'\mid s,a)$ describes action outcomes, the observation model $O(o\mid s',a)$ describes the likelihood of an observation, and $R(s,a)$ and $\gamma$ define the discounted planning objective. Because the true state is latent, the robot maintains a belief $b_t(s)=\Pr(s_t=s\mid h_t)$ over states conditioned on the ordered action--observation history $h_t$.

POMCP~\cite{silver2010monte} provides the online planning backbone used in this work. It performs Monte Carlo tree search over action--observation histories and selects tree actions using
\begin{equation}~\label{eq:ucb}
a^*=\arg\max_a\left[V(ha)+c\sqrt{\frac{\log N(h)}{1+N(ha)}}\right],
\end{equation}
where $V(ha)$ is the estimated action value, $N(h)$ and $N(ha)$ are visit counts, and $c$ controls exploration. The standard planner chooses robot actions; the feedback policy introduced in the next section operates after an action and observation to refine the symbolic belief used by subsequent planning. Thus, external queries are an information-acquisition layer around the planner rather than action choices delegated to the external source.

\subsection{Feedback Decision Problem}~\label{preliminaries:POMCP}
At execution step $t$, the robot has a belief over symbolic successor states after performing $a_t$ and receiving $o_{t+1}$. The proposed policy makes two decisions over this belief:
\[
\begin{aligned}
\textsc{When}(\mathcal{B}_{t+1})&\in\{\texttt{query},\texttt{continue}\},\\
\textsc{What}(\mathcal{B}_{t+1})&\in\mathcal{F}.
\end{aligned}
\]
The first decision should avoid acting on consequential unresolved ambiguity without querying at every opportunity. Conditional on querying, the second decision should select a fact whose answer most effectively separates the remaining state hypotheses. We evaluate these decisions through task success and the number of external queries rather than assuming that either maximal autonomy or maximal information acquisition is always preferable.
